\documentclass[11pt,a4paper]{article}

\usepackage[utf8]{inputenc}
\usepackage[T1]{fontenc}
\usepackage[brazilian]{babel}
\usepackage[margin=2cm]{geometry}
\usepackage{enumitem}
\usepackage{hyperref}
\usepackage{titlesec}
\usepackage{parskip}

\hypersetup{
  colorlinks=true,
  linkcolor=black,
  urlcolor=black,
  pdftitle={Kamille Hartmann Maccagnan - Currículo},
  pdfauthor={Kamille Hartmann Maccagnan}
}

\titleformat{\section}{\large\bfseries}{}{0em}{}[\titlerule]
\titlespacing*{\section}{0pt}{12pt}{6pt}

\pagestyle{empty}
\setlength{\parindent}{0pt}
\setlength{\parskip}{6pt}

\newcommand{\cvheader}[4]{%
  \begin{center}
    {\LARGE\bfseries #1}\\[4pt]
    {\large #2}\\[6pt]
    #3\\[2pt]
    #4
  \end{center}
  \vspace{4pt}
}

\newcommand{\experience}[3]{%
  \textbf{#1} \hfill \textit{#2}\\[4pt]
  #3
  \vspace{8pt}
}

\begin{document}

\cvheader{Kamille Hartmann Maccagnan}%
{Tech Manager | Operações e Finanças | Gestão de Demandas | Liderança e Métricas}%
{Campo Comprido, Curitiba-PR \textbullet{} (41) 9 9928-4424 \textbullet{} \href{mailto:kamillehartmannm@gmail.com}{kamillehartmannm@gmail.com}}%
{\href{https://linkedin.com/in/kamille-hartmann-maccagnan/}{linkedin.com/in/kamille-hartmann-maccagnan}}

\section*{Resumo Profissional}

Profissional com experiência em gestão de demandas, operações de varejo, acompanhamento de entregas e interface entre negócio, finanças operacionais e tecnologia. Atuação na liderança de treinamentos e capacitação, padronização de processos, definição e exposição de indicadores e apoio à tomada de decisão em ambientes com PDV, TEF, POS e fluxos de pagamento integrados a adquirentes como REDE e Cielo. Vivência em homologação funcional do sistema de vendas, investigação de incidentes, priorização de demandas e remoção de impedimentos operacionais. Experiência em PMO, governança de dados e acompanhamento de portfólio de projetos, com análise de riscos, relatórios executivos e alinhamento com stakeholders em ambiente multinacional. Conhecimento em metodologias ágeis (Scrum e Kanban), facilitação de alinhamentos, melhoria contínua e comunicação entre áreas técnicas e de negócio. Pós-graduação em Liderança e Gestão em Tecnologia e inglês avançado.

\section*{Experiência Profissional}

\experience{Anjuss -- Analista de Suporte de TI}{07/2025 -- Atual}{%
Atuação na liderança de treinamentos e onboarding de novos funcionários, com padronização de processos, disseminação de conhecimento e apoio ao desenvolvimento de equipes operacionais em ecossistema de varejo. Responsável pela organização, análise e priorização de demandas em ambiente de operações e finanças de loja, com interface direta com PDVs, TEF e POS integrados principalmente à adquirente REDE, com vivência em ambiente Cielo. Participação em testes em ambiente de homologação do sistema de vendas, com validação de fluxos de pagamento e integração Cielo antes da liberação para as lojas, além de investigação de incidentes em frente de caixa e acionamento de adquirentes para tratativa de falhas. Elaboração de documentação funcional estruturada, incluindo POPs para usuários e materiais técnicos para T.I., garantindo clareza de fluxos, regras e critérios de execução. Desenvolvimento de dashboards em Power BI para monitoramento de indicadores operacionais e reportes de acompanhamento, apoiando visibilidade de entregas, diagnóstico de falhas, priorização de correções e alinhamento entre necessidades de negócio e suporte técnico. Atuação na facilitação de comunicação entre áreas e remoção de impedimentos que impactam a fluidez da operação.}

\experience{HORSE do Brasil -- Estagiária de TI (IS/IT LATAM)}{07/2024 -- 07/2025}{%
Atuação em PMO, governança de dados, suporte e interface com áreas operacionais em ambiente multinacional. Acompanhamento de portfólio de projetos com organização de cronogramas, indicadores de desempenho e visibilidade de entregas, apoiando a quebra de iniciativas em etapas acompanháveis e a validação de entregas junto a stakeholders de negócio, fornecedores e times técnicos. Elaboração de relatórios executivos e materiais de acompanhamento para apresentação de resultados, coleta de informações e apoio à tomada de decisão, com participação em reuniões de alinhamento conduzidas em até três idiomas simultaneamente. Atuação na documentação de processos, análise de riscos, investigação de desvios operacionais, estimativa de impactos e evolução de práticas de gestão, com foco em melhoria contínua, mitigação de incidentes e sustentação operacional.}

\experience{Restaurante Familiar -- Analista de Dados e Suporte Técnico}{01/2023 -- 07/2024}{%
Atuação na análise de necessidades de gestão do negócio e no diagnóstico de problemas que afetavam informações operacionais e financeiras. Mapeamento e estruturação de fluxos de coleta e organização de dados, com criação da base de dados do zero em Excel para viabilizar dashboards em Power BI e definição de indicadores de desempenho. Responsável pela padronização de informações, identificação de inconsistências e proposição de melhorias em processos operacionais com impacto em resultados financeiros. Atuação no alinhamento entre áreas operacionais e administrativas, traduzindo demandas do negócio em estruturas de dados e relatórios com critérios claros de acompanhamento e validação.}

\section*{Projetos e Conquistas}

Estruturação de KPIs e dashboards executivos para monitoramento de indicadores de desempenho e suporte à tomada de decisão. Participação em iniciativas de melhoria contínua e otimização de processos. 1º lugar no Transformation Day Renault 2023 com solução orientada a dados e melhoria de processos. Participação no Hackathon Bosch 2023 com foco em resolução de problemas de negócio com uso de dados.

\section*{Habilidades Técnicas}

\textbf{Liderança e Gestão de Pessoas:} Liderança de treinamentos e onboarding, disseminação de conhecimento, facilitação de alinhamentos, interface com stakeholders, comunicação executiva, pós-graduação em Liderança e Gestão em Tecnologia.

\textbf{Metodologias Ágeis e Entregas:} Scrum, Kanban, Boas Práticas PMOBOK, PMO, acompanhamento de entregas, análise de riscos, melhoria contínua, remoção de impedimentos operacionais, priorização de demandas.

\textbf{Operações, Finanças e Pagamentos:} PDV, TEF, POS, REDE, Cielo, homologação funcional, fluxo transacional em loja, investigação de incidentes, sustentação operacional, contato com adquirentes.

\textbf{Indicadores e Métricas:} KPIs, OKRs, dashboards executivos, relatórios gerenciais, Power BI (DAX, Power Query), Excel, SQL, monitoramento de performance, governança de dados.

\textbf{Ferramentas:} Monday, Trello, ServiceNow, Spotfire, Pacote Office.

\textbf{Idiomas:} Inglês avançado \textbullet{} Espanhol básico \textbullet{} Coreano básico.

\section*{Formação Acadêmica}

\textbf{PUCPR} -- Graduação em Gestão da Tecnologia da Informação \hfill Previsão de conclusão: 06/2026\\
Curitiba, PR

\textbf{Escola Conquer} -- Pós-graduação em Liderança e Gestão em Tecnologia \hfill Conclusão: 2024\\
Curitiba, PR

\textbf{Universidade Tuiuti do Paraná} -- Graduação em Medicina Veterinária \hfill Conclusão: 06/2019\\
Curitiba, PR

\textbf{Qualittas} -- Pós-graduação em Clínica Médica e Cirúrgica de Animais Exóticos e Pets Não Convencionais \hfill Conclusão: 12/2022\\
Curitiba, PR

\end{document}
